\documentclass[11pt,a4paper]{article}
\usepackage[utf8]{inputenc}
\usepackage[T1]{fontenc}
\usepackage{lmodern}
\usepackage[french]{babel}
\usepackage{amsmath,amssymb,mathtools}
\usepackage{icomma}
\usepackage{xcolor}
\usepackage{tikz}
\usepackage{pgfplots}
\usepackage[most]{tcolorbox}
\usepackage{enumitem}
\usepackage{fancyhdr}
\usepackage[hidelinks]{hyperref}
\usepackage{titlesec}
\usepackage[a4paper,margin=2.2cm,headheight=26pt,headsep=10pt]{geometry}
\usetikzlibrary{arrows.meta,positioning,calc,intersections,angles,quotes}
\usepgfplotslibrary{fillbetween}
\pgfplotsset{compat=1.18}
\definecolor{primary}{RGB}{30,75,145}
\definecolor{primarydark}{RGB}{20,50,100}
\definecolor{defcol}{RGB}{0,114,178}
\definecolor{thmcol}{RGB}{0,135,110}
\definecolor{excol}{RGB}{0,130,85}
\definecolor{remarkcol}{RGB}{120,70,180}
\definecolor{attncol}{RGB}{200,40,55}
\definecolor{methcol}{RGB}{85,85,100}
\definecolor{areacol}{RGB}{120,160,230}
\definecolor{areacol2}{RGB}{230,150,80}
\newcommand{\dd}{\mathrm{d}}
\newcommand{\R}{\mathbb{R}}
\newcommand{\ua}{\,\text{u.a.}}
\newcommand{\tg}{\tan}\newcommand{\cotg}{\cot}
\tcbset{boxbase/.style={enhanced,breakable,boxrule=0.9pt,arc=2.5pt,left=3mm,right=3mm,top=2.3mm,bottom=2.3mm,fonttitle=\bfseries,coltitle=white,toptitle=0.6mm,bottomtitle=0.6mm}}
\newtcolorbox{methode}[1][]{boxbase,colback=methcol!7,colframe=methcol,sharp corners,title={\textsc{Comment faire}\ifx\relax#1\relax\relax\else\textmd{ : \itshape #1}\fi}}
\newtcolorbox{exemple}[1][]{boxbase,colback=excol!5,colframe=excol,title={\textsc{Exemple}\ifx\relax#1\relax\relax\else\textmd{ : \itshape #1}\fi}}
\newtcolorbox{idee}[1][]{boxbase,colback=defcol!5,colframe=defcol,title={\textsc{Idée clé}\ifx\relax#1\relax\relax\else\textmd{ : \itshape #1}\fi}}
\newtcolorbox{remarque}[1][]{boxbase,colback=remarkcol!6,colframe=remarkcol,title={\textsc{Remarque}\ifx\relax#1\relax\relax\else\textmd{ : \itshape #1}\fi}}
\newtcolorbox{attention}[1][]{boxbase,colback=attncol!6,colframe=attncol,title={\textsc{Attention}\ifx\relax#1\relax\relax\else\textmd{ : \itshape #1}\fi}}
\newtcolorbox{exo}[1][]{boxbase,colback={primary!4},colframe=primary,title={\textsc{Exercice #1}}}
\newtcolorbox{corr}[1][]{boxbase,colback={thmcol!5},colframe=thmcol,title={\textsc{Corrigé #1}}}
\tikzset{axe/.style={->,>=Stealth,black!65,line width=0.7pt},courbe/.style={primary,line width=1.3pt},courbe2/.style={areacol2!90!black,line width=1.3pt},dvert/.style={gray,dashed,line width=0.7pt}}
\pagestyle{fancy}\fancyhf{}
\fancyhead[L]{\small\textcolor{primarydark}{\textbf{Fiche 7} — Intégrale d'une fonction continue}}
\fancyhead[R]{\small\textcolor{primarydark}{\textsc{Thème 4 — Intégrale définie}}}
\fancyfoot[C]{\small\thepage}
\renewcommand{\headrulewidth}{0.8pt}
\titleformat{\section}{\large\bfseries\color{primarydark}}{\thesection}{0.6em}{}

\begin{document}

\section*{Niveau Difficile — Combiner les propriétés, comparer, découper}

\begin{center}
\itshape\small Ici on raisonne : établir un signe ou un encadrement \emph{sans tout calculer},
combiner plusieurs propriétés, découper avec Chasles, et mener une intégration par parties bien
guidée.
\end{center}

\medskip

\begin{exo}[1 — comparer et encadrer sans (presque) calculer]
\begin{enumerate}[label=\alph*),leftmargin=1.8em]
  \item Sans les calculer, comparer $\displaystyle\int_0^1 x^2\,\dd x$ et $\displaystyle\int_0^1 x\,\dd x$.
        Justifier par une propriété, puis vérifier en calculant les deux intégrales.
  \item Déterminer le \textbf{signe} de $\displaystyle\int_1^2 \ln x\,\dd x$ sans calculer sa valeur.
  \item Établir l'encadrement $\displaystyle 1\leqslant\int_0^1 e^{x^2}\,\dd x\leqslant e$.
        \emph{Indication} : encadrer $e^{x^2}$ sur $[0,1]$.
\end{enumerate}
\end{exo}

\begin{exo}[2 — combiner Chasles, inversion et linéarité]
Soient $f$ et $g$ continues sur $[0,4]$ telles que
\[
\int_0^1 f(x)\,\dd x=3,\qquad \int_1^4 f(x)\,\dd x=-2,\qquad \int_0^4 g(x)\,\dd x=5.
\]
Calculer :
\begin{enumerate}[label=\alph*),leftmargin=1.8em]
  \item $\displaystyle\int_0^4 f(x)\,\dd x$ ;
  \item $\displaystyle\int_4^0 f(x)\,\dd x$ ;
  \item $\displaystyle\int_0^4\bigl(2f(x)+3g(x)\bigr)\,\dd x$ ;
  \item $\displaystyle\int_0^4\bigl(f(x)-g(x)\bigr)\,\dd x$.
\end{enumerate}
\end{exo}

\begin{exo}[3 — intégrale définie par une IPP guidée]
On veut calculer $\displaystyle E=\int_1^{e} x\ln x\,\dd x$.
\begin{enumerate}[label=\alph*),leftmargin=1.8em]
  \item Choisir $U$ (à dériver) et $V'$ (à intégrer) selon la règle de priorité, et donner $U'$ et $V$.
  \item Appliquer la formule $\displaystyle\int_a^b U\,\dd V=\bigl[UV\bigr]_a^b-\int_a^b V\,\dd U$.
  \item Achever le calcul et donner la valeur exacte de $E$, puis une valeur approchée.
\end{enumerate}
\end{exo}

\begin{exo}[4 — quand aire algébrique et aire géométrique diffèrent]
On pose $f(x)=x-1$ sur $[0,2]$.
\begin{enumerate}[label=\alph*),leftmargin=1.8em]
  \item Calculer $\displaystyle\int_0^2 (x-1)\,\dd x$. Commenter le signe du résultat.
  \item En découpant à la racine de $f$ (relation de Chasles), calculer
        $\displaystyle\int_0^2 |x-1|\,\dd x$.
  \item Vérifier sur cet exemple l'inégalité de la moyenne
        $\Bigl|\displaystyle\int_0^2 f\Bigr|\leqslant\displaystyle\int_0^2 |f|$, et interpréter
        géométriquement l'écart entre les deux membres.
\end{enumerate}
\end{exo}

\end{document}
